\pdfminorversion=7
\documentclass[11pt,a4paper]{article}

\usepackage[T1]{fontenc}
\usepackage[utf8]{inputenc}
\usepackage{lmodern}
\usepackage{microtype}
\usepackage[margin=1in]{geometry}
\usepackage{mathtools,amssymb,bm}
\usepackage{graphicx}
\usepackage{enumitem}
\usepackage[authoryear,round]{natbib}
\usepackage[hidelinks]{hyperref}

\hypersetup{
  pdftitle={Refractive Gravity: An Ontological Picture and the Foundations of Mathematical Formalism},
  pdfauthor={George Lotishvili}
}

\graphicspath{{figures/}}
\setcounter{tocdepth}{2}
\setlist{itemsep=0.35em,topsep=0.5em}

\newcommand{\RefG}{\mathrm{RefG}}

\title{\textbf{Refractive Gravity: An Ontological Picture and the Foundations of Mathematical Formalism}}
\author{George Lotishvili\\
\small Associate Professor, East European University\\
\small Tbilisi, Georgia\\
\small \href{mailto:g_lotishvili@eeu.edu.ge}{\texttt{g\_lotishvili@eeu.edu.ge}}}
\date{2026}

\begin{document}

\maketitle

\begin{center}
\textbf{Zenodo DOI:} \href{https://doi.org/10.5281/zenodo.21325652}{10.5281/zenodo.21325652}
\end{center}

% Article body
\begin{abstract}
This paper presents an ontological account of Refractive Gravity (RefG) that brings several lines of scientific work and inquiry together within a single physical picture. RefG is a research program under stepwise development. The central aim of the text is to offer the reader an intuitive picture of the fundamental organization of the Universe; the formalism and completed calculations are available in the public scientific archive on Zenodo.

The current status of the ideas and results discussed in this paper is given in the final section, ``Current Status of the Research Program.''
\end{abstract}

\tableofcontents

\vfill
\begin{center}

\includegraphics[width=0.34\textwidth]{figures/logo.pdf}
\end{center}
\vfill

\clearpage

\section*{Introduction}
\addcontentsline{toc}{section}{Introduction}

Modern physics resembles an intricate puzzle: its individual pieces have been calculated with high precision, tested experimentally, and refined mathematically, yet they still do not fit together into a single physical picture. The Standard Model of particle physics and the geometric description of gravity do not fully fit within one language; inside black holes, classical geometry reaches a singularity; in galactic dynamics, rotation curves and the regularities of modified Newtonian dynamics (MOND) are explained in standard cosmology by dark-matter halos, while the accelerated expansion of the Universe is described in terms of dark energy. Quantum randomness, the nature of time, the hierarchy of particle masses, the fundamental constants, and the relative weakness of gravity belong to the same puzzle. Each has been studied in detail, but a possible common ontological foundation has yet to emerge.

The history of physics offers many unifying visions devised to address these problems. For Albert Einstein, the path to unification lay in a deeper geometry of spacetime. John Archibald Wheeler sought a fundamental picture of the Universe at the boundary between geometry and information, in the principle of ``It from Bit.'' Andrei Sakharov interpreted gravity as an effect arising from quantum fluctuations of the vacuum; Ted Jacobson derived Einstein's equation from the thermodynamic laws of local horizons; and, for Erik Verlinde, gravity is an entropic, holographic phenomenon. Roger Penrose, in turn, treats the global structure of spacetime, cosmology, and quantum physics as a single problem of twistor and cyclic geometry. These approaches differ in their mechanisms, but they ask the same question: can disparate phenomena be described from a single foundation?

This paper assembles ideas developed over many years into a unified logical structure and presents the ontological framework of \textbf{Refractive Gravity (RefG)}.

The standard language of modern physics generally begins with a ready-made apparatus of measurement: coordinates, a time parameter, and fields are taken as given from the outset. RefG begins at the boundary that precedes them: before we can speak of space, time, or mass, nature must contain a physically distinguishable, persistent difference. What physical condition gives these concepts the form of measurable quantities?

In 1905, Einstein abandoned absolute time and defined simultaneity operationally through clocks and light signals \citep{einstein1905}. RefG poses this operational question at a prior level: what physical condition creates the clock and signal on which those operations themselves depend?

To answer this question, the paper rests on three fundamental principles:

\begin{itemize}
\item \textbf{First:} Nature contains only one universal underlying medium. What we call fields, forces, and particles are different phase, pressure, shear, rotational, or topological modes of this single substance.
\item \textbf{Second:} The stable unit of matter is the \emph{oscillon}---a locally confined, self-sustaining resonant form of the underlying medium. Its intense internal dynamics leave a persistent pressure-deficit imprint in the surrounding medium.
\item \textbf{Third:} Gravity is the refractive response of the underlying medium. What we perceive as attraction is the wave-like motion of an object through an effective geometry modified by the pressure deficit.
\end{itemize}

The Michelson--Morley experiment of 1887 detected no ether wind of the kind expected from the Earth's motion and placed stringent constraints on a stationary, matter-independent luminiferous ether \citep{michelson1887}. The underlying medium of RefG belongs to a different ontological picture: space, time, the measuring apparatus, and the observer are all modes of this one foundation. RefG associates uniform motion with the symmetric wave and phase transfer of form, which leaves no dissipative imprint on the corresponding Lorentz-covariant branch. This picture revives and extends Lord Kelvin's deep intuition: matter is itself a form of the medium, and the underlying medium is the substance from which the Universe is made.

The present text develops the ontological picture of RefG, in which the computational results collected in the public Zenodo archive acquire an intuitive physical meaning. It connects the underlying medium with measurability, the pressure deficit associated with mass with cosmological horizons, and the origin of particles with the topology of quantum numbers. Medium-based models of unification have appeared before, including Winterberg's Planck aether \citep{winterberg2002}. Within this line of thought, RefG places particular emphasis on the operational origin of measurability and on the mechanism of population-wide resonant locking. The following is a non-exhaustive list of the principal contributions of this work:

\begin{enumerate}
\item \textbf{Population-locked time:} The local rate of time, \textbf{$\Omega(x)$}, is generated by mutual locking within the particle population itself, \textbf{$\nu_i=n_i\,\Omega(x)$}; it does not rely on an externally prescribed metronome. Time dilation is a collective reconfiguration of this local network. This population-based definition of the temporal rate is a distinctive element of RefG.
\item \textbf{Two-channel accounting:} By definition, gravitational mass is a readout channel through the medium. The oscillon's internal state is retained as a physically separate quantity, and RefG distinguishes the two axiomatically.
\item \textbf{Separation of extensive and intensive properties:} On the selected first-order biconformal branch, operational quantities scale according to a \textbf{$p/p^2$} law; RefG relates these powers to the ontological distinction between nodes and connections.
\item \textbf{The Koide $C_3$ phase structure:} In the three-phase map of the charged leptons, $\theta=h/9$ is the reduced coordinate of an order-nine return cycle. The first nontrivial return of the oriented charged frame, $h=2$, corresponds to $\theta=2/9$, while equalization of the internal channels corresponds to $A_K=\sqrt{2}$.
\item \textbf{Node-based interpretation of the fine-structure constant ($\alpha$):} Within RefG, $\alpha$ is most plausibly a dimensionless invariant ratio between the electric and magnetic responses of the underlying node network. Operationally, it measures the relative strength with which the same network supports elementary electric flux and a transverse photon oscillation.
\item \textbf{Quarantine of process time:} A clear boundary is drawn between two languages of time---coordinate time and process time. Process time remains an interpretive layer and does not enter the actual calculations of the primary cosmological observables.
\item \textbf{Two readouts of one operator:} The unified spectral map of RefG associates short modes with the particle readout and long modes with the cosmic-web readout.
\end{enumerate}

Priority is assessed separately for each link in this chain. Medium-based programs of unification have precursors, and $p/p^2$ scaling also appears in the polarizable-vacuum tradition \citep{dicke1957,puthoff2002}. RefG derives these powers from the ontology of nodes and connections and joins them to two-channel mass accounting.

The three-phase parametrization of the Koide relation and its geometric reading continue the work of Brannen and Foot, while a different holonomy-based derivation of $2/9$ is due to Shulga \citep{brannen2010,foot1994,shulga2026}. RefG links this structure to the $\theta$-blindness argument, the order-nine return cycle, and equalization of the internal channels.

The relaxation-based reading of dark energy and the knot-based argument for three-dimensionality likewise build on established lines of thought \citep{klinkhamer2008,whitrow1955,berera2017}. The distinctive backbone of this paper is their connection to the birth of measurability, the population-wide $\Omega(x)$, two-channel accounting, and the two scale-dependent readouts of a single spectral operator.

This picture associates dark energy with the remnant of the historical relaxation of the underlying medium, and three-dimensionality with knot closure and topological stability.

Modern information, symbolic, and software tools enable the systematic computational testing of the intuitive ideas developed in this research.

The text invites the reader to follow this chain of reasoning and see how one physical picture connects phenomena that otherwise appear separate. The scientific value of the framework will be determined by its further formalization and observational testing.

\section{Ontological Foundation}

\subsection{A Single Underlying Medium: Substance, Emptiness, and the Dynamics of Existence}

Modern fundamental physics describes the Universe as a collection of many independent fields. In the Standard Model, each fundamental field has its own degrees of freedom and rules of interaction. Refractive Gravity (RefG) replaces this picture with a single fundamental ontological principle: nature contains only one universal underlying medium. What we perceive as different fields and forces are, in fact, different topological, phase, or rotational modes of this one substance. Such a monistic approach has deep roots in the history of physics. As early as 1870, William Clifford proposed that matter was nothing more than small, moving ``hills'' of spatial curvature \citep{clifford1876}. Yet there is an essential difference between Clifford's vision and the RefG framework: Clifford's carrier is purely geometric, whereas our underlying medium supports pressure, shear, and resonant channels in addition to its geometric response.

When the Universe is conceived as a single medium, even a minimal picture of its dynamics requires two asymptotic limits. At one pole lies the absolute maximum of substance---a state in which pressure, mass--energy potential, and operational density approach their limiting values. At the other lies emptiness, or the zero limit, where physical interaction vanishes. This ``emptiness'' is a purely logical limit at which physical measurability falls to zero. Objective reality is the continuous, bounded oscillation between these two poles. The minimal active units of the underlying medium---``existence nodes''---emerge within this continuous dynamics.

The ontological unity of the underlying medium is accompanied by structural richness. It has several independent channels: it can be compressed, undergo torsion, form vortex flow, or transmit a wave excitation. Water offers a useful analogy: the same fluid can simultaneously support hydrodynamic pressure, bulk flow, and surface tension, although these properties are not directly reducible to one another. Likewise, the universal underlying medium combines phase, pressure, shear, and topological responses. This multichannel character supplies the resources from which electric charge, spin, and the complex particle spectrum are formed.

The inactive state of the underlying medium, denoted formally by $\varphi=0$, is also one of its states. The medium is ubiquitous; what varies is its local degree of excitation, the density of distinguishable nodes, and the strength of the connections between them. In this language, a drop in pressure, or emptiness, is an increase in operationally unmeasurable intervals.

Here ``pressure'' denotes the effective pressure--energy variable of RefG. The condition $\varphi=0$ denotes the background state, while excitation relative to it can develop in either direction: $\varphi>0$ is a state of elevated pressure--energy, whereas $\varphi<0$ is the deficit branch. The gravitational readout chain in this paper is built on the profile of the deficit branch. Classical hydrodynamic pressure is an intuitive analogue of this variable. In the external readout, a change in $\varphi$ appears as a change in the rate of time, the coordinate propagation of signals, operational size, and effective mass.

\subsection{The Origin of Measurability and Operational Space}

As long as the underlying medium remains in a state of absolute, inactive uniformity, its physical measurement is impossible in principle. In such a state, the concepts of ``distance'' and ``time'' do not exist, because there is no second physically distinguishable point for comparison. Physics proper begins with the medium's self-differentiation, or self-distinction: once a local asymmetry---a physically distinguishable node, a phase displacement, or a pressure deficit---arises within the medium, it no longer remains a trivial, indistinguishable background.

The interval that arises between any two distinguishable nodes is the physical seed of space; the delay in their connection is the seed of time; and the energetic tension of differentiation is the ontological seed of mass. External gravitational mass is defined later, through the readout of the deficit profile created by this tension.

Physical meaning requires comparison. Three elements appear at the ontological origin: a \textbf{node}, a locally distinguishable state of the underlying medium; an \textbf{imprint}, the change that the node leaves in the surrounding medium; and a \textbf{relation}, the possibility of comparing one state with another. From this triad unfolds the operational network that, in macroscopic language, we call space, time, and matter.

This view is central to understanding the nature of the observer. We, like our measuring instruments, are complex composite structures built from the same nodal network. The approach strongly echoes the classical tradition of operationalism, in particular Bridgman's principle that a physical concept is defined by the operations used to measure it \citep{bridgman1927}. Refractive Gravity goes further: it seeks the fundamental ontological condition that makes those measuring operations themselves possible. In this account, the measurable Universe is not matter placed within an abstract ``space''---it is the continuous process by which the underlying medium distinguishes itself. Our operational space arises from the imprint of this distinguishability. Beyond that imprint, the universal medium may continue to exist, but it appears to us as a physical object only where it leaves a measurable, operational consequence.

The ontological distance between nodes of the underlying medium is not measured directly. The observer is formed from modes of this same medium and has no measuring device outside it. What can be measured is the operational imprint of the interval between nodes. Operational space appears when the relation between at least two nodes becomes physically distinguishable. Below this threshold, the difference remains a local internal state of the medium; once the threshold is crossed, it is read through changes in the rate of time, operational size, mass, and signal propagation.

RefG posits a minimum threshold of operational distinguishability, denoted by $\ell_*$. In this picture, no independent measurable relation has yet arisen below the scale $\ell_*$; at $\ell_*$, the first stable, distinguishable record appears; and a larger interval between nodes is read through the operational signatures of pressure, the rate of time, and signal propagation. The physical motivation for this assumption is that an independent record at an indefinitely small scale could not retain its stability. The Planck length is a natural unit constructed from the constants $c$, $G$, and $\hbar$, whereas $\ell_*$ is the hypothetical dynamical threshold at which a stable record arises. Whether the two coincide depends on the microdynamics of the underlying medium.

It follows from this fundamental principle that the ontological origin of the Universe is not a multiplicity of points. If an operational interval between nodes has not yet formed, then two physically distinguishable nodes do not yet exist, and neither does a spatial relation between them. Such a state cannot even be called a geometric ``single point,'' because the very concept of a point presupposes that space has already formed. It is more precise to call it a completely indistinguishable state of the underlying medium. Multiplicity appears in nature only when a difference arising within the medium leaves a stable physical imprint and the relation between nodes becomes operationally distinguishable; from this differentiation emerge the familiar structures of space, time, and the local rhythm of motion.

It is precisely at this fundamental boundary that the primary physical content of gravity comes into view. When the local connection between nodes of the underlying medium weakens, the external observer reads this ontological state as lower pressure, slower clocks, and altered conditions of signal propagation. A macroscopic observer perceives this coordinated ensemble of physical properties as a ``gravitational field.'' The following chapters develop this conceptual sketch into a quantitative refractive chain.

\subsection{The Local Collective Rhythm and Resonant Locking}

The shared rhythm of local time is not a preset metronome frequency imposed on the underlying medium from outside. It arises naturally when stable oscillons---the wave-like forms of matter---join in a common resonant network and their intrinsic frequencies settle into fixed, stable ratios to one another. In this collective state, the rate of time is not determined by an external background; it is created directly by the internal coordination and mutual locking of the local network itself.

At a deeper ontological level, the underlying medium is neither a tone nor a metronome. It is the instrument: a three-dimensional vibrating foundation that sets the boundaries, provides the resources, and determines the vibrationally admissible spectrum. A bell offers a vivid analogy. Its sound is not produced by isolated notes; the modes of a single body resonate together. Likewise, the frequencies of elementary particles are established through mutual balance and coordination within the resources of the medium. The universal medium determines which modes can exist, while the particle modes balance one another within that admissible space and form a stable, locked configuration.

Denoting the shared local rhythm by $\Omega(x)$, the frequencies of stable oscillons may be written as

\begin{equation}
\nu_i(x)=n_i\,\Omega(x),
\label{eq:population-locking}
\end{equation}

Here $n_i$ is the fixed relative factor of the mode in question. In general, it is not an integer, so the particle frequencies do not reduce to simple harmonic overtones. They more closely resemble the mildly anharmonic spectrum of a complex three-dimensional volume resonator. We shall return to the three-phase structure of the charged leptons in later chapters.

The relation $\nu_i=n_i\Omega$ leaves each mode with its own frequency while tying them all to one local rhythm. The ratios $n_i/n_j$ are therefore preserved under resonant locking. A rhythm incompatible with the network loses phase coherence and cannot persist as a stable particle mode.

In a low-pressure region, the collective rhythm $\Omega(x)$ of the entire local resonant network changes. In a dense medium, coupling is strong and $\Omega$ is higher; in a rarefied medium, coupling weakens and $\Omega$ falls. Slower transfer between nodes, a lower energetic state of the medium, and a decrease in $\Omega$ are three readouts of one state; each enters the quantitative chain only once.

When $\Omega(x)$ changes, every local clock and particle process scales with it. A local observer reads the rate of their own clock as normal. Under this common scaling, dimensionless ratios---including $\alpha$ and particle-mass ratios---remain invariant because they express internal structural ratios of the network rather than its absolute rhythm.

The classical antecedent of this relational picture is Mach's principle \citep{mach1883}. RefG gives it a local form: the shared rhythm is set directly by the local resonant network. The relational description is preserved, and physical action remains local.

In the RefG map, the fine-structure constant is read as a dimensionless invariant ratio between the electric and magnetic responses of the underlying node network.

Now consider a wave imprint propagating from one pressure regime of the underlying medium into another. Its observed frequency is calculated from the effective metric. In the ontological language of RefG, the comparison between the local rhythms of the source and receiver is read as a phase translation, while cosmological redshift is defined by the metric relation $1+z=a_0/a_{\rm em}$. The imprint does not carry a frozen absolute frequency along its path: in each region, its phase relation is translated to the local $\Omega(x)$. A compatible imprint locks into the new resonant network, while an incompatible one dephases or scatters; its energy is redistributed among compatible pressure--energy channels of the underlying medium.

Ultimately, within the RefG framework, the rate of time is the ontological state of the local resonant network itself. When a change in the pressure of the underlying medium transposes the network by a common factor, all compatible frequencies change together, while their dimensionless ratios remain invariant. An independent universal clock has neither an operational role nor any additional measurable consequence in this description: every real comparison takes place within the same physical network. This principle has a clear test---if, under a common transposition, any dimensionless ratio changes with $\Omega(x)$ in a different way, the postulate of a shared population-wide rhythm will be ruled out.

\subsection{The Oscillon, the Pressure Deficit, and the Origin of Mass}

What we conventionally call a ``particle''---whether an electron or a proton---is treated in the present framework as an oscillon. An oscillon is a locally confined, stable wave structure in the universal underlying medium. It is not an isolated rigid body independent of the medium; on the contrary, it is a highly organized, self-sustaining state of that medium. The concept has an interesting historical antecedent in Lord Kelvin's model of vortex atoms \citep{thomson1867}. Although that particular nineteenth-century program has passed into history, its conceptual core---the particle as a self-sustaining dynamical form of a medium---is the deep intuition that the oscillon picture revives and carries forward in modern resonant and topological language.

When an oscillon vibrates or rotates intensely, or stores energy within its form, the static pressure of the surrounding medium falls. Bernoulli's picture makes this direction vivid \citep{bernoulli1738}: an increase in dynamical intensity is accompanied by a decrease in static pressure. On the selected weak, static branch of RefG, this analogy corresponds to a pressure--energy balance identity.

By existing as an active form, an oscillon generates a pressure-deficit profile in the underlying medium. In this map, the deficit is the direct external profile of gravitational mass: its integral readout gives the system's gravitational charge. The more intense the oscillon's internal dynamics, the deeper the deficit and the stronger the refractive gradient. The oscillon's internal reserve is accounted for separately from this external readout.

This picture has a well-known historical parallel in John Archibald Wheeler's program of ``mass without mass,'' in which a geon is a massive configuration held together by the energy of its own field \citep{wheeler1955}. Wheeler's geons gradually radiate their energy away; RefG associates the stability of an oscillon with resonant locking and topological protection.

The term \emph{effective refractive mass} denotes the external readout of an individual oscillon in a given background. Its internal state---its topological structure and internal energy---is invariant in the local description, while the external gravitational readout changes with the background. The asymptotic mass of the complete system belongs to a separate integral readout channel and does not follow automatically from the $m_{\rm eff}$ of an individual oscillon. With this distinction, the origin of mass is read as the formation of the medium's refractive response.

At this fundamental level, gravity is the refractive deflection of the wave form of one oscillon within the pressure deficit created by another. A slack string provides an intuitive image of the gravitational well: the locally reduced pressure--energy state of the underlying medium weakens the conditions of propagation. The deficit profile shapes the effective geometry of the medium, and the particle's wave form moves refractively through that geometry.

An oscillon generates a pressure deficit, and the deficit in turn changes its effective mass, operational size, and coupling to the medium. RefG reads this two-way feedback as a mechanism of dynamical equilibrium: the material form changes the medium, while the altered medium determines the conditions of the form's external physical manifestation.

\section{Gravity as Refraction}

\subsection{Three Concurrent Operational Changes}

As noted in the preceding chapter, the interval between nodes of the underlying medium is not measured directly; only its physical imprint is operationally accessible. Around an oscillon, this imprint appears as a local reduction in the pressure state of the medium: physically distinguishable nodes are more sparsely distributed, and the informational connection between them is weaker. For an external macroscopic observer, this single ontological state manifests itself through three simultaneous, coordinated changes:

\noindent\textbf{First: a reduction in the coordinate speed of light.} Within a pressure-deficit region, nodes are more sparsely distributed and their interaction is weaker, so light and other wave signals propagate more slowly in the coordinate description. The locally measured speed of light remains unchanged; what changes is the propagation rate read in the coordinate language of a distant observer.

\noindent\textbf{Second: a reduction in the operational size of the oscillon.} The fundamental topological form of the oscillon remains unchanged, while in a rarefied underlying medium it is filled by fewer interacting nodes. In the external description, this is read directly as a smaller operational size of the oscillon and, correspondingly, as a modified effective-mass readout.

\noindent\textbf{Third: a reduction in the rate of the local resonant clock.} A clock is a real system built from oscillons, and its rhythm depends on local nodal connections. In a rarefied medium, those connections weaken, and the clock slows in the coordinate time of a distant observer. The local clock itself continues to count its own rhythm as normal.

A double-walled vacuum flask offers a vivid analogy: fewer carrier molecules mean less frequent collisions and slower heat transfer. In the RefG map, this corresponds to the relation between the strength of nodal coupling and the coordinate rate of a clock.

These three changes are three operational readouts of a reduced pressure state. Their quantitative dependence on the change in the medium is not the same.

\begin{figure}[htbp]
\centering
\includegraphics[width=0.92\linewidth]{figure_3.pdf}
\caption{Operational readout of an oscillon in a pressure-deficit region.}
\label{fig:oscillon-pressure-deficit}
\end{figure}

\noindent\textbf{(A)} In a high-pressure underlying medium, the oscillon occupies a dense nodal configuration; the connection between nodes is strong, and local processes proceed at a high rate.

\noindent\textbf{(B)} In a pressure-deficit region, the nodes of the underlying medium are more sparsely distributed; the operational separation between nodes increases, the rate of coupling falls, and local time slows.

\noindent\textbf{(C)} In the external readout, an oscillon within the pressure deficit is read with a slower temporal rate, a smaller operational size, and a modified effective mass; the enlarged internal interval does not enter directly as a measurable quantity.

The separation of nodes does not generate a ``second space.'' The interval between nodes is not an independent physical object; within this framework, ordinary ``emptiness'' is read through the pressure variable. The operational manifestation of a rarefied underlying medium is a local reduction in its pressure state.

\subsection{The Quadratic Relation and the Factor of Two}

The three operational changes do not scale with the same power. In this map, the operational size and effective mass of an individual oscillon belong to the node-count channel, whereas the coordinate speed of light belongs to the channel governing the strength of inter-node coupling. The first readout is extensive; the second is intensive, meaning that it does not depend on the size of the system.

Writing the pressure factor as $p$, the scaling takes the form

\begin{equation}
\frac{L_{\rm oper}}{L_0}
=\frac{m_{\rm eff}}{m_0}
=\frac{\Omega}{\Omega_0}
=p,
\qquad
\frac{c_{\rm coord}}{c_0}=p^2,
\qquad
p=e^{\varphi/2}.
\label{eq:operational-scaling}
\end{equation}

Similar powers also appear in polarizable-vacuum models \citep{dicke1957,puthoff2002}. RefG relates them to the ontological distinction between nodes and connections.

In formal terms, the external macroscopic readout of an individual oscillon is related to the pressure-deficit profile of the underlying medium by a specific exponential half-factor:

\begin{equation}
m_{\rm eff}=m_0 e^{\varphi/2}.
\label{eq:effective-mass}
\end{equation}

On the deficit branch, $\varphi<0$, and hence $0<p<1$. This expression describes the external readout of an individual oscillon placed in a given background: as $\varphi$ decreases, its operational size and $m_{\rm eff}$ decrease. The total gravitational charge of the system that generates the deficit is defined separately through an integral readout. This relation follows from the metric filters and the selected first-order biconformal branch of the theory.

The following four quantities must be kept distinct:

\begin{enumerate}
\item the oscillon's internal state and its own energy;
\item the effective mass $m_{\rm eff}$ of an individual oscillon in a given background;
\item the profile of the deficit source generated in the underlying medium;
\item the ADM/Komar-type asymptotic gravitational charge of the complete system \citep{arnowitt1962,komar1959}.
\end{enumerate}

These quantities form two nested balances:

\begin{equation}
M_{\rm proper}=M_{\rm external}+M_{\rm hidden},
\qquad
M_{\rm external}=M_{\rm active}+M_{\rm dressing}.
\end{equation}

The first equality separates the internal reserve from the external readout; the second resolves the external readout itself into an active core contribution and a dressing carried into the surrounding medium. The half-factor for an individual oscillon, the three-dimensional volumetric filter, and the boundary charge are different readout maps. In asymptotic accounting, the readout of the active core and the dressing of the surrounding medium are gathered into a single boundary charge, excluding any double counting of the same imprint.

On the selected static branch, the projection of the deficit profile is read as an effective four-component source for the metric. A similar distinction exists in standard gravity: the local energy of a test body and the integrated mass of the entire system are not the same quantity.

A vivid analogy brings these complex, self-consistent dynamics into view. Imagine a colored, ice-like crystal immersed in clear water. The crystal is a locally bound form of the same water, just as an oscillon is a local resonance of the common underlying medium. As the crystal dissolves, its color passes into the water: the surrounding medium gradually becomes colored, while the crystal itself shrinks. Energy does not disappear---like color dissolving into the water, it passes from a locally bound state into a broader state of the medium.

As more of the underlying medium becomes ``colored'' by this imprint, its local pressure state falls, and the externally read mass, operational size, time, and distance change. This is a relaxation process: the more of its imprint the oscillon transfers to the surrounding medium, the less capacity it retains to transfer any further imprint. The process therefore slows, and the deepening of the deficit is limited by its own accumulated imprint.

Against an external measurement standard, this evolution proceeds toward minimum distinguishability---possibly toward a scale of Planck order---and at the limit the object appears to vanish from the external readout. Local measuring instruments participate in the same relaxation regime, so this limit arises differently in the internal readout. The mass deficit and the energy transferred to the underlying medium are two independent readings of a single process.

This scaling has the biconformal metric representation

\begin{equation}
ds^2=e^\varphi c^2dt^2-e^{-\varphi}d\boldsymbol{\sigma}^2,
\qquad
\varphi=-\frac{2U}{c^2},
\label{eq:bi-conformal-metric}
\end{equation}

where $U>0$ is the potential of the selected exponential branch and, at leading order, agrees with the magnitude of the Newtonian potential. In the weak-field regime,

\begin{equation}
g_{tt}=1-\frac{2U}{c^2}+\frac{2U^2}{c^4}+\cdots,
\qquad
g_{ij}=-\left(1+\frac{2U}{c^2}+\cdots\right)\delta_{ij}.
\label{eq:weak-field-expansion}
\end{equation}

On the selected first post-Newtonian (1PN) weak-field branch, these expressions give the parametrized post-Newtonian (PPN) coefficients $\beta=1$ and $\gamma=1$.

The duration of one complete clock cycle is given by the ratio of operational size to coordinate speed. If the size is halved while the coordinate speed falls by a factor of four, the duration of the cycle doubles---the clock runs at half its original rate.

On the same branch, the first post-Newtonian results agree with the standard values for light deflection, Shapiro delay, gravitational redshift, and perihelion precession \citep{bertotti2003,will2014}. The Cassini radio-link test constrains $\gamma-1$ to a precision of about $2\times10^{-5}$. Historically, the additional factor of two supplied by the spatial component completes Einstein's 1911 half-result to the full general-relativistic deflection \citep{einstein1916,shapiro1964}. The higher-order Solar-System branch is selected by a separate field equation and does not follow automatically from this exponential representation.

\subsection{Refractive Attraction, Inertia, and the Absence of Friction}

An oscillon, as a wave structure, exists within a gradient of pressure and coordinate propagation speed in the underlying medium. Phase propagates faster on the denser side and more slowly on the more rarefied side, so the wavefront is refracted toward the low-pressure region. In the weak, static, large-scale limit, the deficit created by an object shapes the effective geometry, while wave propagation through that geometry manifests itself as gravitational motion.

In RefG, the fall of an apple has the following picture: the wave forms of its atoms bend toward the region of lower pressure and slower coordinate propagation created by the Earth. Familiar attraction is the refractive response of the medium.

In formal terms, the gravitational source is the pressure deficit projected into the underlying medium. It is the component of the pressure--energy balance that an external observer reads as attraction. Microscopically, it is the stress profile imprinted on the medium by the oscillon's internal dynamics.

A gravitational well is a three-dimensional region of reduced pressure in the underlying medium. The mass of a bound system is lower than the sum of its freely separated constituents by the binding energy.

GW150914 provides an observational example: gravitational waves carried away energy equivalent to approximately three solar masses \citep{abbott2016}. In general relativity, this energy balance is the standard physics of binding energy. RefG gives it its own reading in the language of the deficit profile: the merged node acquires a deeper external profile, while the remaining energy of the reconfiguration leaves as a transverse--traceless (TT) wave. A complete dynamical calculation determines the quantitative relation between the source profile and the radiated flux.

On the static compact branch, the projected deficit source is read as an effective violation of the null energy condition (NEC). In the RefG map, this stress belongs to the confinement channel of the underlying medium.

RefG associates inertia with the deformation of the deficit imprint during acceleration. The pressure profile ahead of and behind the oscillon becomes asymmetric, and this asymmetry produces a reaction opposite to the acceleration. Under braking, the picture is reversed as in a mirror.

On a localized, finite-energy, Lorentz-covariant branch of the oscillon, when the Laue boundary condition is satisfied, the leading inertial coefficient corresponds to the total Noether energy \citep{noether1918,laue1911}. Under the same condition, inertial and gravitational masses are read together at leading order. The MICROSCOPE experiment tests this universality at the level of $10^{-15}$ \citep{touboul2022}.

On the branch of a uniformly moving Lorentz-covariant solution, the wave pattern is transmitted symmetrically and leaves no dissipative imprint. RefG associates frictionless motion with this regime of phase transmission.

Motion here is a wave-like relay of form. Its limiting speed is determined by the maximum rate of phase transfer between nodes of the underlying medium. The finite speed of light arises from the same operational network.

\subsection{Local Unobservability and the Principle of Relativity}

A local observer in a gravitational field reads their own temporal rate and spatial scale as normal, because their clock, ruler, and body are all made from the same underlying medium. As the environment changes, the entire local system is reconfigured by a common factor. A change sufficiently slow in comparison with the internal periods is adiabatic: it preserves the mode and excites no unwanted transitions. Under this common scaling, dimensionless physical ratios, including $\alpha$, remain invariant.

This picture is subject to firm experimental bounds. Gravitational redshift has been tested across many regimes, from the Pound--Rebka experiment to Gravity Probe A and the Galileo tests \citep{pound1960,vessot1980,delva2018,herrmann2018}. Quasar spectra and atomic clocks independently constrain the stability of dimensionless ratios; in cosmological comparisons, the typical bound on $|\Delta\alpha/\alpha|$ is of order $10^{-6}$. RefG reads these facts as physical manifestations of local self-calibration.

The null result of the Michelson--Morley experiment corresponds to common-mode self-calibration, whereas the LIGO signal is a differential tensor response. When the underlying medium changes uniformly throughout a region, the ruler, clock, and conditions of light propagation are all reconfigured in the common mode, and the local readout self-calibrates. Modern tests of Lorentz symmetry constrain directional differences in the speed of light to the order of $10^{-17}$. A gravitational wave is a time-dependent, quadrupolar differential deformation: it acts with different signs along different axes and therefore leaves a measurable imprint in a detector.

The FitzGerald--Lorentz length compensation is a historical precursor of this logic \citep{fitzgerald1889,lorentz1892}. RefG associates the compensation with the fact that measuring instruments are themselves wave excitations of the underlying medium. In the local Minkowski limit, the tensor branch is luminal and produces a transverse--traceless (TT) detector response. Comparing the gravitational and electromagnetic signals from GW170817 constrains $|v_{\rm gw}-c|/c$ to the order of $10^{-15}$ \citep{abbott2017}, while the orbital decay of binary pulsars tests the normalization of quadrupolar energy loss at the level of $10^{-3}$--$10^{-4}$ \citep{weisberg2016,kramer2021}. The scalar breathing branch describes an additional volumetric deformation of the underlying medium; its contribution is constrained by the GWTC-3 tests of gravity \citep{ligovirgokagra2022}.

The clearest operational invariant in this ontological picture is the difference between the proper times accumulated by two clocks that follow different trajectories and later reunite. The numbers recorded on their displays are objective consequences of their histories through different gravitational regimes. Locally measured rates, sizes, and frequencies remain relational readouts of processes unfolding within a common environment.

\subsection{Historical Context and the Place of RefG}

The refractive reading of gravity has clear precursors in the history of physics. As early as 1911--1912, before the final formulation of general relativity, Albert Einstein related the coordinate speed of light to the Newtonian gravitational potential through a relation of the form $c=c_0(1+\Phi_N/c^2)$ \citep{einstein1911,einstein1912}. The tradition of medium-based optics has still deeper roots: in 1851, Hippolyte Fizeau showed that moving water partially ``drags'' the propagation of light \citep{fizeau1851}, while Walter Gordon related the propagation of light through a dielectric medium to the geodesics of an effective metric \citep{gordon1923}. This established a formal bridge between the optics of media and gravitational geometry.

In 1957, Robert Dicke recast gravity in the language of a variable refractive index of the vacuum \citep{dicke1957}. The starting point of modern analogue gravity is William Unruh's result that acoustic excitations in a moving fluid propagate in an effective metric analogous to that of a black-hole horizon \citep{unruh1981}. Volovik's condensed-medium program and the line developed by Barceló, Liberati, and Visser showed how an effective metric can emerge from the density and dynamics of a background medium \citep{barcelo2011}. This line of research also received a striking laboratory continuation: in 2019, a thermal analogue of Hawking radiation was observed in a Bose--Einstein condensate \citep{munozdenova2019}.

Alongside this physical-medium tradition runs the thermodynamic branch of emergent gravity. Andrei Sakharov interpreted gravity as an effect generated by quantum fluctuations of the vacuum \citep{sakharov1967}, while Ted Jacobson showed that Einstein's equation can be obtained as an equation of state from the thermodynamic conditions of local horizons \citep{jacobson1995}. This branch is related to RefG in its view that spacetime geometry is not fundamental; in RefG, effective geometry is generated by a concrete pressure--refractive medium.

Refractive Gravity adds three fundamental new elements to this rich historical line:

\begin{itemize}
\item \textbf{First---an ontological foundation:} A single universal underlying medium is introduced that provides the domain through which light propagates and also constitutes matter itself---oscillons, measuring instruments, and the observer. Through this ontological step, the approach moves from an auxiliary mathematical ``analogy'' to an independent framework for physical research.
\item \textbf{Second---a shared local resonant rhythm:} The collective rhythm of oscillons turns the local unobservability expressed by relativity into an internal dynamical consequence of the system.
\item \textbf{Third---a universal pressure-deficit channel:} RefG uses the pressure deficit as a common language of source and readout. This language connects mass, the rate of time, light deflection, and the local tensor response as different readouts of a single pressure--energy balance.
\end{itemize}

Within this historical map, the closest conceptual comparison for RefG is Erik Verlinde's entropic gravity \citep{verlinde2011}. In Verlinde's model, Newton's law arises from the entropic balance of information on a holographic screen. RefG shares the idea that gravity is emergent and takes the pressure deficit of the underlying medium as its source. The two mechanisms anticipate different observational outcomes and are tested independently.

\section{Three Scales and Cosmological Horizons}

\subsection{The Local Scale: Compact Objects, Internal State, and External Readout}

The universal pressure chain developed step by step in the preceding chapters appears in nature on three entirely different astrophysical scales: locally, as the pressure deficit formed around an oscillon; galactically, as a coherent flow of the universal underlying medium; and cosmologically, as a historical, global change in the state of that medium. In fact, all three phenomena are different dynamical regimes of one and the same underlying medium. We begin at the local scale, where this chain assumes its sharpest and most extreme form---in compact objects.

In the extreme-gravity regime, an object's \emph{internal state} becomes sharply distinct from its \emph{externally read mass}. RefG's compact-object map distinguishes the internal state of accumulated matter from the mass read at a distance. A distant observer reads that state through the deficit profile, the rate of time, and the operational size. The internal geometry, horizon, and causal structure are determined by the complete solution.

The physical meaning of this distinction is simple: the object's own material content does not disappear inside. One unit remains one, and two remain two; compression and strong binding do not nullify the internal inventory. What changes is its external readout---the deficit gradient, the rate of time, the conditions under which light can escape, and the operational size. This is precisely the two-channel accounting of RefG: the internal inventory and the externally read active mass are different physical quantities.

RefG associates the mass difference of a merged system with changes in binding and in the deficit profile. Its quantitative relation to the radiated tensor flux is determined by the complete dynamical solution. In Chapter II, GW150914 provided an observational example of this energy map.

The metric determines the difference between internal and external scales. On the compact branch of RefG, that difference can be large. The exact internal scale, the existence of a horizon, and the global causal structure are determined by the complete solution.

Viewed from outside, such an object may appear small, dark, and almost frozen, even while the proper length of the same region is much greater. The locally measured speed of light is unchanged in both descriptions; what differs is the temporal and spatial readout received by a distant observer. The ``frozen'' appearance is a regime of external readout, whereas the question of a horizon is settled by the complete causal structure. The historical intuition of a dark star already appears in Michell's work of 1784 \citep{michell1784}.

In this language, a singularity marks the limit of a description that cannot unite the internal inventory, a finite core, the pressure deficit, and the external readout within a single self-consistent system.

Penrose's singularity theorem rests on several assumptions, including the relevant energy and causal conditions \citep{penrose1965}. On the compact branch of RefG, the effective source violates the null energy condition (NEC), so the theorem's conclusion does not automatically extend to this branch. A regular interior regime is determined by the complete field solution, its boundary matching, and its stability.

The idea of regular compact objects has also been developed in the modern literature. A well-known example is the ``gravastar'' of Mazur and Mottola, a condensate object without a singular core \citep{mazur2004}. Such alternatives may be distinguished by their shadow size, oscillation spectrum, and echo signals \citep{cardoso2019}. RefG is structurally related to this family and ties the interior regime to the pressure--energy balance of the universal underlying medium.

RefG associates the feedback between mass and deficit with a finite equilibrium state. The field equations determine its existence and stability. In this picture, a compact object is a heavily loaded regime of the underlying medium whose internal state and remotely read mass are registered through different channels.

\subsection{The Galactic Scale: Vortical Flow and MOND}

\paragraph{Observational law.} Galaxy rotation curves do not fall off as the distribution of visible baryonic matter alone would suggest. At the same time, baryonic mass and rotation speed obey a tight relation---the baryonic Tully--Fisher relation. Historical landmarks are the rotation curve of Andromeda \citep{rubin1970} and the Tully--Fisher relation itself \citep{tully1977}. In the modern radial acceleration relation, 2,693 measurements from 153 galaxies lie on a single common curve \citep{mcgaugh2016}.

The Lambda cold dark matter model ($\Lambda$CDM) describes these phenomena with dark-matter halos. MOND instead singles out a characteristic acceleration scale, $a_0$, at which departures from the Newtonian expectation set in \citep{milgrom1983}. Tensor--vector--scalar theory (TeVeS) and its newer generalizations use additional vector and scalar fields to provide a relativistic description of this phenomenology \citep{bekenstein2004,skordis2021}. RefG associates $a_0$ with the coherence scale of the flow of the underlying medium.

\paragraph{The vortical mechanism of RefG.} The structure behind this map is a large-scale \textbf{vortical, or swirling, flow} of the underlying medium. An additional pressure deficit forms at the center of this ordered rotational regime, and baryonic matter gathers in the resulting well. The flattening of rotation curves and the Tully--Fisher relation are read as the combined refractive effect of a giant vortex in the underlying medium and the imprints of baryonic oscillons.

The closest modern parallel to this picture is superfluid dark matter, in which a MOND-like force is associated with phonon excitations of a condensate phase \citep{berezhiani2015}. There, superfluidity is a phase of an additional particle component; in RefG, coherent flow is a regime of the single underlying medium itself.

\paragraph{Morphology.} A galaxy's visible form may reflect the flow regime of the underlying medium. Ordered, laminar rotation favors disk and spiral structures; a more chaotic or isotropic flow favors dispersion-supported systems. These are two endpoints of a continuous spectrum, since some elliptical galaxies also retain substantial rotation.

\paragraph{Clusters.} In galaxy clusters, modifying the acceleration law alone often fails to close the dynamical mass budget. Historically, the missing-mass problem first came to light in observations of the Coma Cluster \citep{zwicky1933}. RefG reads the cluster mass budget through three distinct contributions: the global nodal deficit of the long-wavelength resonant network; the galaxies and their local oscillon tails; and cluster-scale vortical flow together with the memory left after a merger. All three channels enter one common pressure-deficit accounting, with each contribution counted once. In Chapter IV, particles and the cosmic web will appear as the two ends of this same spectrum.

Two discriminating observational signatures emerge from this picture. A continuous field can also produce a vortex, so swirling flow by itself does not distinguish whether the underlying medium is discrete. The signature of nodal microstructure is quantized circulation together with a compatible dispersion relation.

\paragraph{First---quantization of circulation.} If the flow of the underlying medium has a single-valued condensate phase, the expected unit of circulation is $h/m$, where $m$ is the mass of the effective flow mode. For an extremely small $m$, the imprints of individual vortices may appear in weak gravitational lensing or in disk kinematics. Similar signatures are sought in the wave-like halos of ultralight dark matter \citep{hu2000,hui2021}.

\paragraph{Second---the frozen imprint of the Bullet Cluster.} In the Bullet Cluster, the centers inferred from weak gravitational lensing are offset from the X-ray plasma at a significance of approximately $8\sigma$ \citep{clowe2006}. RefG uses this separation as a quantitative benchmark for a model of frozen memory. The lifetime of that memory must be determined by a coupled model of transport in the underlying medium and gas hydrodynamics; the discriminating data are the lensing map and the spatial distributions of the gas and galaxies.

The resulting test map is unified: it spans the galactic data for the SPARC radial acceleration relation (RAR) \citep{lelli2016,mcgaugh2016}, galaxy clusters, Bullet-Cluster-type mergers, the cosmic microwave background (CMB), and large-scale structure. On a branch without dark matter, the vortical source must also replace the cosmological contribution of cold dark matter.

Holding fixed the metric branch, matter content, recombination history, and initial conditions, RefG leaves the linear sources of the primary CMB unchanged; this is a test of algebraic consistency. In a cosmology without dark matter, the role of cold dark matter passes to the vortical and nodal channels, and the true test is whether this replacement reproduces the full likelihood of the CMB and large-scale-structure data.

\subsection{The Cosmological Scale: Relaxation of the Underlying Medium and Two Languages of Time}

At the cosmological scale, the external metric description remains unchanged: the Universe expands, distances between galaxies increase, and light undergoes cosmological redshift. Viewed from within, RefG reads this history as the relaxation of the underlying medium. The long-range imprints of oscillons accumulate into a common background and alter the pressure state of the medium. In the process-time description, earlier epochs correspond to a larger internal budget of cycles for structure formation.

Photons follow null geodesics of the effective metric, and cosmological redshift obeys the standard scale-factor relation, $1+z=a_0/a_{\rm em}$. In the ontological language of RefG, this corresponds to the phase translation described in Chapter I: an imprint arriving from an earlier state of the underlying medium is read against the local resonant rate in each new epoch, while dimensionless ratios are preserved.

The ``age of the Universe''---approximately 13.8 billion years---is metric time calibrated by present-day clocks \citep{planck2020}. Before the first stable material modes formed, no operational clock had been defined. Process time provides an additional description of the internal history of the underlying medium.

The idea of two time scales in cosmology has historical antecedents. Milne described the same Universe in terms of a kinematic time $t$ and a dynamical time $\tau$ \citep{milne1935,milne1948}. Wetterich linked the languages of expansion and changing material standards within a specific field model \citep{wetterich2013}. In RefG, process time describes the budget of the medium's internal cycles.

In a self-similar process description, each epoch measures the past with its own local clocks. A future epoch corresponding to 15.8 billion years on our clock could, under its own process calibration, still read its entire past as approximately 13.8 billion years. The metric age continues to increase in the usual way; self-similarity applies only to the rule by which internal cycles are counted.

A gradually fading vibration of a string makes these two languages vivid. An internal clock counts time by the number of oscillations, and its unit gradually lengthens; a fixed external clock records the same decay at a different rate. One physical history is read on two different measurement grids.

In early JWST data, massive red galaxies appeared approximately $500$--$700$ million years after the Big Bang \citep{labbe2023}. A substantial part of this tension is resolved by star-formation efficiency, revised mass estimates, and the contribution of active galactic nuclei. The process budget concerns only the remaining moderate discrepancy.

Cosmological expansion is expressed in RefG at three descriptive levels:

\begin{enumerate}
\item \textbf{In the external metric language,} distances in space increase.
\item \textbf{In the material language,} clocks and atomic standards from different epochs occupy different local states of the underlying medium.
\item \textbf{In the language of the underlying medium,} the mean density of nodes and the connections between them change globally over cosmic history.
\end{enumerate}

\begin{quote}
These are three levels on which the same history can be described. The external metric is the primary language of observational calculations.
\end{quote}

The limit of the medium's relaxation is defined thermodynamically: diminishing gradients mark an approach to maximum entropy. RefG reads this process as an \textbf{asymptotic} evolution in which equilibrium remains a limiting state. Complete equalization would erase the very contrast between substance and emptiness whose smoothing drives the process; relaxation therefore approaches the limit without terminating at any finite time.

RefG associates the macroscopic arrow of time with the direction in which the underlying medium relaxes. The dissolving-crystal picture introduced in Chapter II makes this intuition vivid: color dispersed through water does not spontaneously gather back into a solid crystal. RefG reads this statistical irreversibility as the macroscopic direction of time.

In this interpretive map, the thermodynamic, cosmological, and radiative arrows share a common statistical source in the relaxation of the underlying medium. Macroscopic irreversibility reflects the one-way approach to equilibrium \citep{planck1899}, while the arrow of time is tied to the medium's initial nonequilibrium state.

The late Universe is expanding at an accelerated rate. Standard cosmology attributes this phenomenon to $\Lambda$ or dark energy. RefG reads it as an energetic residue left by the relaxation of the underlying medium. The stress of this residue gives rise to negative pressure, while its slow variation yields an almost constant cosmological term. RefG relates the numerical value of $\Lambda$ to the parameters of the relaxation equations.

Within the ontology of a single medium, the vacuum-energy problem becomes a question of the equilibrium stress--energy tensor of the underlying medium and its gravitational readout. The conventional sum of zero-point oscillations of independently quantized fields, whose estimate exceeds the observed value by approximately a factor of $10^{120}$, is no longer the starting form of the accounting in this ontology: all channels enter the common equilibrium of a single self-sustained medium. RefG reads the small cosmological residue as a departure left over from this equilibrium. Related mechanisms have been studied in condensed-matter physics and in $q$-theory \citep{volovik2003,klinkhamer2008}. The same underlying medium gives rise to both gravitational refraction and late-time relaxation. The discriminating quantities in this reading are $\rho(a)$, $w(a)$, and $\Lambda$; their observational coordinates are baryon acoustic oscillations (BAO), Type Ia supernovae, and the CMB \citep{abdulkarim2025,brout2022,planck2020}.

\paragraph{The process-time quarantine.} Process time is an internal interpretive scale. The expansion history, the CMB, Big Bang nucleosynthesis (BBN), photon redshift, atomic clocks, and stellar nuclear processes are calculated in metric time. The spectral evolution of supernovae follows the observed $(1+z)$ time dilation \citep{blondin2008}. The process scale does not enter the formulae for these primary observational calculations.

\subsection{Quantum Randomness and Two-Level Time}

The twofold nature of time also offers an interesting interpretive possibility in quantum physics. The place and moment of a measurable event are registered in oscillon time---that is, clock time. RefG associates the selection of an outcome with the internal level of the underlying medium; this level is not measured directly by clock time, because the clock itself is an oscillon structure produced by the same medium.

Violations of Bell inequalities rule out local hidden-variable models that explain outcomes through locally preassigned parameters \citep{bell1964,hensen2015}. At the same time, the no-signaling condition preserves relativistic causality: distant correlations cannot be used to transmit controllable information faster than light.

RefG reads this coexistence as \textbf{atemporal internal coordination}. In the language of the underlying medium, it does not describe a signal traveling through space; in the language of clocks, it appears only as a correlation between distant outcomes.

Wheeler's delayed choice is read in the same language \citep{wheeler1978}. In a single-photon realization, the measurement arrangement chosen at a late stage determines which complementary outcome appears: interference in a closed interferometer, or path information in an open one \citep{jacques2007}. RefG's two-level time describes these facts without a retrocausal signal.

This picture has influential relatives in modern theoretical physics: the timeless fundamental level of the Wheeler--DeWitt equation; the Page--Wootters mechanism, in which time emerges from internal correlations among subsystems; and Bohm's pilot-wave picture, with its nonlocal hidden coordination \citep{dewitt1967,pagewootters1983,bohm1952}. RefG enters this family with its own ontology of the underlying medium.

This quantum picture has two strict quantitative criteria:

\begin{enumerate}
\item the \textbf{Born probability distribution}---the macroscopic statistical readout of internal coordination \citep{born1926};
\item the \textbf{no-signaling condition}---internal coordination cannot become a transmissible signal that violates relativistic causality.
\end{enumerate}

\section{Oscillons and the Origin of Matter}

\subsection{Many Patterns in One Ocean}

For a vivid picture, imagine an unbounded, continuous ocean with waves rising within it. Most of them form, propagate, and soon disappear, but a few settle into stable, confined patterns: one may be ring-shaped, with the flow circulating continuously as a vortex; another, a stable spherical pulsation; and a third, a spiral wave.

\clearpage

\begin{figure}[htbp]
\centering
\includegraphics[width=0.95\linewidth]{figures/ai-render-16799135.jpg}

\caption{Resonant modes of the underlying medium. In the intuitive map of RefG, stable wave patterns in a single medium are read as different particle modes: a spherical pulsation, a ring-shaped vortex, and a spiral wave. Here a particle is not a separate substance; it is a locally confined resonant pattern of the underlying medium.}
\label{fig:resonant-modes}
\end{figure}

The ocean picture conveys the central proposition of this chapter: the electron, proton, photon, and quark are different resonant patterns of one universal underlying medium. The medium itself is what oscillates; the particle is a stable form of that oscillation. A historical precursor is the Skyrme model, in which nucleons are read as topological solitons of a nonlinear meson field \citep{skyrme1961}. RefG takes this intuition one step further: every such form is a mode of the same underlying medium.

The classical image of a small, solid sphere gives way here to a membrane-like picture: the vibrational amplitude of the underlying medium extends through space, while a localized event in a detector is the resonant contact between this vibration and the detector's internal oscillon modes. In this way, RefG gives the spatial amplitude and phase of the quantum state an ontological foundation in the underlying medium.

The electron cloud is the amplitude itself, not a hidden path rapidly traversed by a small body. The extended state carries a real phase structure through space, while the point of registration arises through local resonant contact with the detector.

Single-electron interference brings two facts together: the state's amplitude depends on the geometry of both slits, yet each registration appears as a single point on the screen. Tonomura's group directly recorded the gradual build-up of this pattern \citep{tonomura1989}. Likewise, a stationary electronic state in an atom does not radiate as a classical orbit would. RefG reads these results as the conjunction of an extended wave in the underlying medium and localized registration.

The width of a wave packet and the particle's internal structure are different things. Scattering measures the internal form factor of charge and energy; this is how, for example, the proton radius is read. The packet's spatial extent does not enter this observable as an internal size. In high-energy scattering, the electron retains a pointlike electromagnetic form factor down to scales of approximately $10^{-18}\,\mathrm m$ \citep{navas2024}. A direct discriminating test of the lepton-oscillon map is whether the form factor derived from the charge profile reproduces this same pointlike behavior over the accessible momentum range.

\subsection{Localized Confinement: Massive and Massless Modes}

In RefG's particle map, the principal distinguishing feature between massive and massless modes is \textbf{localized confinement}.

In this map, the energy of a massive oscillon is confined within a localized form. When vortical circulation or stable pulsation closes a wave pattern within a finite region of space, it creates a pressure-deficit profile. Long-lived localized pulsations are known in nonlinear field theory as pulsons and oscillons \citep{bogolyubsky1976,gleiser1994,copeland1995}. RefG associates their stability with population-wide resonant locking. In the external readout, this confined energy appears as inertia and gravitational mass.

Ordinary field-theoretic oscillons are long-lived, but they gradually decay through radiation. RefG associates the suppression of radiative leakage, and the selection of long-lived modes, with mutual locking across the population.

A massless carrier mode has neither a rest frame of its own nor a stable localized core; it propagates through space at the limiting speed. RefG interprets the photon as a propagating phase mode of this kind.

Helicity and confined local rotation are different things. The photon's helicity $\pm1$ describes the helical structure of phase along the direction of propagation; it does not create vortical circulation confined to one place. A massive mode has a stable localized profile, whereas a massless carrier has a propagating dynamical phase mode.

A photon's energy and momentum enter the gravitational balance. Free radiation leaves no persistent, localized deficit profile of the kind associated with a massive oscillon, but its energy contributes in full to the source of the complete system.

The same rule extends to composite particles. Strong-interaction energy confined within a hadron is not read as a separate external imprint, but it contributes to the hadron's total energy and gravitational mass.

Local form, spin, charge, color, and polarization retain the internal distinctions among particles; the gravitational readout combines them in a common total pressure-deficit channel. Binding energy, the energy of free radiation, and all of the particle's remaining energy content therefore enter the gravitational source of the complete system.

\subsection{Self-Forming Resonators and the Two-Level Filter}

A wave equation permits an effectively infinite variety of forms, yet nature presents us with a sparse, well-defined spectrum of stable particles. Why does nature select this particular set from infinitely many possibilities? RefG's answer rests on a \textbf{two-level filter}.

The first level is a kinematic and population filter. A handpan provides the intuition. The construction of a guitar predetermines the resonant length of its string; a three-dimensional oscillon, by contrast, must ``forge'' its own resonant cavity within the medium. In a stable handpan note, the fundamental, the octave, and the fifth above the octave form a triplet in the ratio $1:2:3$. Here this triplet is a structural rhythm: a self-forming resonator brings several mutually attuned modes together within one body. In the RefG map, an oscillon's deficit imprint changes the local geometry, and the altered geometry in turn selects the viable modes. A mode must be able both to propagate through the underlying medium and to coexist with the existing particle spectrum. Stability is the result of compatibility across the entire local population.

The second level is an energetic-stability filter. A form that passes the first filter but has an open kinematic decay channel transitions into lower-energy modes. Hadronic resonances of the $\Delta$, $\rho$, and $\omega$ type are familiar examples. Electron decay has not been observed, and charge conservation strongly protects the electron as the lightest charged particle. The observed stability of the proton requires the appropriate baryonic symmetry protection; RefG's particle dynamics must use that protection to account for the lower bound on its lifetime.

Typical lifetimes of hadronic resonances are of order $10^{-22}$--$10^{-24}$ seconds: they appear in the spectrum, but cannot remain as long-lived members and decay into already allowed lower-energy modes. The energetic filter therefore does not erase the spectrum; it redistributes energy among stable channels.

The mathematical language of this filter is that of \textbf{attractors} in nonlinear dynamics---stable states that a system approaches asymptotically as it evolves. The energy window selects which compatible modes survive as long-lived particle states.

A stable particle is a self-consistent equilibrium: the wave pattern creates a deficit in the medium, and through feedback that deficit selects and reinforces the pattern itself. In nonlinear fields, Q-balls are protected as localized configurations by a conserved internal charge \citep{coleman1985}. In RefG, protection also arises from phase locking across the entire population; it does not require the introduction of a new individual charge. The gravitational feedback between mass and deficit is read as the microscopic continuation of this same equilibrium.

The underlying medium is the common vibrational channel, while phase locking occurs directly among the oscillons. They synchronize with the existing spectrum through the fixed resonant ratios described in Chapter I.

This mutual locking requires neither a master clock nor an external conductor. Huygens's pendulums suspended from a common beam, phase-locked laser modes, and the Kuramoto model provide familiar examples of spontaneous synchronization \citep{huygens1665,kuramoto1975}. A historical relative is the bootstrap program in particle physics, in which the spectrum is an ensemble of mutually compatible modes \citep{chew1961}.

A sufficiently slow change in the gravitational field preserves these locks. The change is adiabatic when it proceeds slowly enough relative to the internal periods to avoid exciting unwanted transitions. In a rapidly changing, strongly nonadiabatic background, the mismatch with the surrounding medium may exceed the binding strength and break the lock; the released energy is redistributed among modes compatible with the local configuration.

When a wave imprint enters a different pressure environment, it is read anew under the local operational conditions: its wavelength, process rate, operational size, confinement scale, and externally read mass all change together, while dimensionless ratios and phase harmony are preserved. The particle configuration retunes collectively within the new resource regime. This is the transposition principle described in Chapter I.

\subsection{The Global Resonator: Unity of the Microscopic and Macroscopic Spectra}

The two-level filter tells us \emph{how} the particle spectrum is selected. We now ask: where is it realized? Ernst Chladni's plate provides the intuition \citep{chladni1787}. Sand scattered across a vibrating metal plate gathers along the nodal lines of a standing wave, where the motion is minimal. In the RefG map, this picture corresponds to global modes and their nodes. Every resonator has a spectrum of its own: geometry and boundaries determine which modes are amplified and which fade away.

In RefG's global-resonator picture, the Universe is read as a \textbf{global volumetric resonator}---a three-dimensional underlying medium with its own allowed spectrum. The strong excitation at the cosmological beginning activates this spectrum as a whole. Modes reinforce or cancel one another; their deficit imprints lock into resonance or drift out of phase. A form incompatible with the geometry and background pressure cannot remain a stable particle, and its energy passes into the allowed channels.

In this picture, today's stable particles are the local, short-wavelength modes of the universal spectrum. Where phase ordering did not have time to complete in the early Universe, frozen lines of topological defects may have remained in the underlying medium. We return to this cosmological imprint in Chapter V.

In the same map, dispersion in the underlying medium links frequency to spatial wavelength. Short modes correspond to localized particles; long modes, to the large-scale distribution of matter---the cosmic web.

Galaxies here resemble Chladni's sand: their distribution reveals where the underlying medium has stable nodes, filaments, and vast voids. Galaxy clusters are the deepest nodes in this network, where long-wavelength resonance lowers the background pressure. This connects directly to the picture of clusters developed in Chapter III.

The conclusion of this map is both simple and beautiful: the microscopic particle and the macroscopic cosmic web are two readouts of one universal spectrum---a form bound locally on the inside, and a pattern spread globally on the outside. In standard cosmology, the filaments, sheets, and voids of the cosmic web are described by Zel'dovich's anisotropic collapse \citep{zeldovich1970}. RefG connects this morphology to a resonant ontological language.

RefG's unified spectral map associates short modes with the particle readout and long modes with the background nodal deficit. A joint test of this operator encompasses the localized particle spectrum, the history of structure formation, the mass budget of clusters, and the BAO standard ruler at approximately $100\,h^{-1}\mathrm{Mpc}$ \citep{eisenstein2005}. It is this joint test that distinguishes two readouts of one spectrum from a mere splice of two independent descriptions.

RefG reads the Universe as a finite, boundaryless resonator. Wetterich's model shows that a change of variables can recast expansion as a relation between geometry and material standards \citep{wetterich2013}. In RefG, the compatibility of these languages is tested through the agreement of all dimensionless observables.

In the finite-resonator picture, the size of the underlying medium is fixed, while cosmological history is translated through changes in internal standards and operational readout. The CMB, BBN, and $T\propto(1+z)$ are calculated in the metric language. A possible finite-size imprint at the scale of the horizon would appear in low-multipole power and in cosmological topology.

Present data show no clear discrete global spectrum: the significance of the low-multipole anomalies is modest, while searches for cosmological topology push the possible size of the resonator beyond the observable horizon \citep{planck2016}. The detection of a global cutoff would be a necessary signature of finiteness, though it would not constitute unique evidence; failure to detect one only raises the lower bound on the size.

In the internal language of a fixed resonator, the initial singularity is read as a degenerate initial configuration of local standards. In the same language, the question ``what does the Universe expand into?'' no longer arises: no external vessel is required, because metric expansion is a relation between the geometry of the underlying medium and its own material standards.

\subsection{Discreteness and the \texorpdfstring{$C_3$}{C3} Symmetry of the Generations}

In the Standard Model, the masses of the electron, muon, and tau lepton are empirical parameters. RefG reads them as an internal, order-nine, three-phase ($C_3$) triplet of a single topological defect: three distinct resonant modes of one physical structure.

The parametrization of the charged leptons combines the relation $m\propto\nu^2$, or $\nu\propto\sqrt m$, with a three-phase structure:

\begin{equation}
\nu_k=\mu\left[1+A_K\cos\left(\theta+\frac{2\pi k}{3}\right)\right],
\qquad k=0,1,2.
\label{eq:c3-triplet}
\end{equation}

Any positive triplet can be written in this three-phase form. Its physical content is determined by two dimensionless parameters---the modulation amplitude $A_K$ and the phase $\theta$---and by their dynamical origin.

Here $\mu$ is the common frequency scale, $A_K$ is the amplitude of the three-phase modulation, $\theta$ is the initial phase, and $k$ labels the three generations. In RefG's order-nine return map, the oriented charged branch with $h=2$ corresponds to $\theta=2/9$ radians. The relation $m\propto\nu^2$ reads mass as the integrated energy of a locally confined oscillation.

For $A_K=\sqrt2$, the $C_3$ algebra gives the empirical Koide ratio, $K=2/3$, exactly \citep{koide1982,koide1983}. The three-phase parametrization and its geometric reading are known in the literature: Brannen recast it in the language of a circulant phase \citep{brannen2006}, while Foot related $K=2/3$ to the $45^\circ$ geometry of the $\sqrt m$ vector \citep{foot1994}. Brannen's later Pancharatnam--Berry construction yielded $\pi/12$ and left $2/9$ as a dynamical phase \citep{brannen2010}. In an independent compact family cycle, Shulga obtained $\theta_\ell=-2/9$, with the opposite orientation sign \citep{shulga2026}. The route taken by RefG is different: $2/9$ is the reduced phase of the order-nine return cycle of a single physical defect.

A simple $C_3$-invariant potential cannot select $\theta$. This $\theta$-blindness places the source of the phase beyond the minimum of a simple potential, in framed topology and holonomy. The second structural number, $A_K=\sqrt2$, expresses an equal distribution of energy between the internal singlet and doublet channels; the origins of the two numbers are attributed to independent mechanisms.

The value $A_K=\sqrt2$ is associated with equalization of the internal channels, while $\theta=2/9$ corresponds to the $h=2$ branch of the order-nine cycle. Anchored to the electron pole mass, the exact branch $(A_K,\theta)=(\sqrt2,2/9)$ misses the muon mass by approximately $451.7\sigma$ and the tau mass by $0.61\sigma$. This difference fixes the required magnitude of the radiative bridge between the geometric form and the pole masses \citep{sumino2009}.

This balance also has a simple musical intuition: the charged-lepton $C_3$/Koide triplet resembles the cold, long-lived, slightly anharmonic resonance of a three-dimensional volumetric resonator. The tones are different modes of one body.

\section{Topology and Quantum Numbers}

\subsection{Phase, Orientation, and Topological Quantization}

Up to this point, we have discussed the universal underlying medium primarily in the language of amplitude and local pressure. The full local state of the medium is described by a collection of fields: the phase--clock scalar, the amplitude profile, and the spatial frame each have their own variables. In this article, $\Phi$ is shorthand for this unified underlying medium, not a bare one-component scalar; the full state carries amplitude, phase, internal form, and orientation together.

If $\Phi$ is only a single-valued scalar and the flow is defined directly by $u=\nabla\Phi$, the vorticity vanishes. A topological sector requires a multivalued phase, a defect line, and an internal orientation locked in space. The amplitude vanishes at the defect core, while the phase winds around it an integer number of times. A quantized vortex in a superfluid is a laboratory example of such a structure.

The cosmological counterpart of this structure is the \textbf{cosmic string}---a macroscopic defect line frozen in during a phase transition in the early Universe. Its formation is described by Kibble's causal mechanism \citep{kibble1976}. In the language of RefG, this is the “frozen seam” mentioned in Chapter IV. The vacuum topology of the underlying medium determines whether such lines form, while observations constrain their cosmological density.

Lord Kelvin's work on vortex atoms and vortex motion \citep{thomson1867,thomson1869} led Peter Tait to catalog aether knots and helped give rise to the mathematical theory of knots \citep{tait1877}. RefG carries this historical line forward through pressure--energy balance and population-wide resonant locking.

This historical line also carries a clear lesson. Vortex atoms in an ideal fluid could neither produce a stable particle spectrum nor reproduce the atomic data. In RefG, a knot is the kinematic skeleton of a form; its viability is determined by pressure--energy feedback and population-wide resonant locking.

Phase quantization arises when the field is continuous along a loop, the state returns to the same single-valued state after a complete circuit, and the loop belongs to a nontrivial topological class. Under these conditions, the total change in phase is an integer multiple of $2\pi$.

The Călugăreanu--White--Fuller relation applies to a closed, orientable, framed ribbon \citep{calugareanu1961,white1969,fuller1971,ricca2011}:

\begin{equation}
Lk = Wr + Tw
\label{eq:self-linking}
\end{equation}

$Lk$ is the integer linking number, $Wr$ describes the spatial writhe of the axis, and $Tw$ the internal twist of the frame. The geometrical contributions $Wr$ and $Tw$ may change, but their sum $Lk$ is invariant under the corresponding continuous deformations.

In hydrodynamics, helicity is conserved in an ideal fluid, and Moffatt connected it to the linking number \citep{moffatt1969}. Onsager and Feynman predicted the quantization of circulation in superfluid helium, and Vinen measured the quantum of circulation experimentally \citep{onsager1949,feynman1955,vinen1961}. Integer phase winding is thus part of direct laboratory physics.

In the RefG map, the formula \textbf{“topological locking = charge”} means that electric charge is the external readout of a locked phase--orientation path. The topological number fixes the discrete class, while the field dynamics describe its electromagnetic readout.

The Möbius strip does not lie within the direct domain of this relation; it is described by an orientable double cover. We shall use the spinorial intuition of this double cover in the next subsection.

In RefG's synthetic map, spin-$1/2$ and unit charge are read as two geometrical readouts of a single double-cover source: the half-turn determines the spinorial return of the state, while the orientation of the twist and the integer winding determine the sign of the charge and its discrete class. This common source connects them, although the spinorial transformation and the normalization of electric flux are independent quantitative structures.

This direction also offers a testable indication. Kelvin waves have been directly excited on quantized vortices in superfluid helium, and their dispersion has been measured \citep{minowa2025}; helical axial modes and their dispersion have also been measured directly in a classical water vortex \citep{barckicke2026}. In both cases, a Kelvin wave is a helical disturbance that travels along the vortex axis, although the dispersion coefficients depend on the medium and on the structure of the core. The relevant object of comparison in RefG is therefore not the name but the oscillon's internal dispersion law: the radial mode is linked to a Compton-like mass pulsation, and the helicoidal mode to the radiation channel.

\subsection{Spin and Fermionic Statistics}

A phase twisted by half a turn along the oscillon's internal axis creates a double-cover structure like a Möbius strip. One complete circuit reverses the orientation; a second restores the initial state. This $Z_2$ double cover bears the central signature of a spinorial transformation: a $2\pi$ rotation changes the sign of the state, while a $4\pi$ rotation restores it completely. RefG associates the electron's spin-$1/2$ with this geometry. Changing the topological class requires the amplitude to vanish at the core, the defect to reconnect, or the boundary condition to change.

This geometry also provides an intuition for fermionic exchange. When two identical double-cover modes are exchanged, each passes through the other's twisted phase field, and the exchange path shifts the phase by half a cycle. A full recovery of the Pauli principle requires this shift to translate, in two-particle space, into a minus sign under exchange, an antisymmetric wave function, and zero amplitude when both particles occupy the same one-particle state. An initial version of this vortex-exchange route was formulated in the author's earlier work \citep{lotishvili2026b}.

In a bosonic mode, the exchange phase does not change the sign of the state, so many bosons can occupy the same quantum state. A vivid example is the phase-aligned, or coherent, photonic state of a laser.

This is the topological intuition behind fermionic exchange. In relativistic quantum field theory, the spin--statistics relation is a rigorous theorem based on Lorentz invariance, positivity of energy, and microcausality \citep{pauli1940}. The quantitative criteria for RefG are the satisfaction of these conditions and the antisymmetric construction of Fock space.

\subsection{Charge, the Photon, and the Electromagnetic Channel}

RefG reads electric charge as the external imprint of an oscillon's internal phase twist. The two opposite orientations of the twist correspond to positive and negative charge. Outside the oscillon, this imprint spreads in every direction, and its total flux through a closed surface determines the magnitude of the charge.

In this map, the positron is the topological mirror image of the electron. When a particle and its antiparticle annihilate, their opposite phase windings combine, and the released energy passes into the radiation channel. In a neutral mode, the contributions of opposite phase twists balance, and the net flux through a closed surface vanishes. RefG associates the neutrino with this neutral sector.

Turning the asymmetry angle by the same amount at every point leaves the physical picture unchanged. A measurable imprint arises when this angle varies in space. RefG reads the photon as a transverse, helicoidal wave of that variation: it propagates the phase change of the asymmetry and mediates the electromagnetic interaction.

The topology of the electromagnetic field also appears in classical Maxwell theory. The field lines of Rañada's vacuum solutions form linked loops and are classified by the Hopf index \citep{ranada1989}. RefG connects this mathematical picture to the underlying-medium channel of charge and electromagnetic dynamics.

The sign of the force is determined jointly by the mediator's spin, the type of source, and the sign of the coupling. In the electromagnetic vector channel, like charges repel and opposite charges attract; in the tensor gravitational channel, positive energy produces universal attraction. In three spatial dimensions, the dilution of a spherical flux gives a $1/r^2$ law, while the corresponding potential follows a $1/r$ dependence.

Bjerknes hydrodynamic forces show the limit of this analogy clearly: bodies pulsating in phase in a fluid attract, whereas bodies pulsating out of phase repel \citep{bjerknes1906}. In electrostatics, the signs are distributed in the opposite way. A mechanical resemblance therefore cannot determine the sign of the force; that sign is fixed by the joint structure of the vector mediator, the source, and the coupling.

The fine-structure constant $\alpha$ measures the dimensionless strength of this electromagnetic channel. In RefG, it specifies how the same node network responds, in relative terms, to electric flux and a transverse photon oscillation. This ratio remains invariant under a common rescaling. The quantitative targets of this channel are Maxwell's equations, Gauss's law, charge conservation, $E=\hbar\omega$, bosonic statistics, the creation and annihilation of a free photon, and the mirror self-conjugacy of its two helicity states.

\subsection{Color, Confinement, and Three Generations}

Three-dimensionality acquires a structural role in the RefG particle map. The three spatial directions serve as a geometrical carrier for the internal color triplet. Color remains an internal degree of freedom; the spatial axes provide its geometrical basis.

A structure oscillating along a single axis is topologically incomplete and has open ends. Combining all three axes corresponds to baryonic closure, while joining an axis to an anti-axis corresponds to mesonic closure. An isolated quark cannot form a closed state; RefG associates confinement, the trapping of color, with this geometry. Quarks carry three color states, gluons mediate transitions between those color states, and leptonic modes are color-neutral.

Here $D_3$ denotes the discrete deformation used in the eight-mode count. The resulting eight modes numerically match the eight generators of $SU(3)$ and form a structural map of the color sector. The discriminating criteria for full $SU(3)$ are the algebra of these modes, local symmetry, and their couplings.

The question of three generations is the logical continuation of the $C_3$ triplet discussed in the preceding chapter. In the RefG reading, charged leptons are the ninth-order internal three-phase modes of a single topological defect, while the quark generations carry the same generational architecture into the colored axial sector. At this level, a generation is no longer a replication added to a table from outside; it is a recurring mode of internal topological locking in the underlying medium.

The familiar foundation for holonomy is Berry's geometrical phase: during cyclic evolution, a state accumulates, alongside its dynamical phase, a phase that depends on the geometry of the path it has traversed \citep{berry1984}. RefG connects this language to the return cycle of a single topological defect. Berry's result establishes the existence of holonomy; the specific value $2/9$ is selected by the ninth-order internal return map.

Color and generation are different structures. Color indicates which internal axial direction is active; generation indicates which stable mode of the same resonator has been selected. Their triplicity is a numerical correspondence, not a structural identity.

In RefG, three-dimensionality is tied to the stability of knotted loops. One-dimensional closed loops form nontrivial knots precisely in three spatial dimensions; in higher dimensions, additional directions appear through which they can be untied. This argument singles out three-dimensionality in models where matter is built from such loops \citep{ehrenfest1918,whitrow1955,berera2017}.

Here the same stability filter performs two selections: it selects the particle spectrum from mutually compatible closed forms, and it singles out three spatial dimensions as the setting for their topological protection. In RefG's synthetic picture, the discreteness of matter and the three-dimensionality of space are two consequences of one locking principle.

\subsection{A Geometrical Map of the Standard Model}

\subsubsection{The Electron and Leptons}

In this map, the electron is a complete toroidal structure with a Möbius-type twist, in which all three spatial axes close locally. RefG associates this double-cover geometry with the Dirac value $g=2$. The measured gyromagnetic factor of the electron is $g_e=2.00231930436\ldots$; Schwinger calculated the first term in the deviation, $a_e=\alpha/(2\pi)$, while a modern single-electron experiment measures $g/2$ to a relative precision of approximately $1.3\times10^{-13}$ \citep{schwinger1948,fan2023}. This value of $g_e$ is a strict quantitative target for the topological map.

\subsubsection{Quarks and Color Charge}

In this map, quarks are axial segments with open ends. Their inability to exist in isolation corresponds to confinement. RefG considers two geometrical routes to fractional charge: assembling a unit winding from three minimal substeps, and locking the topological phase in thirds within a confined state. Color and electric charge are different channels; an electric charge of $2/3$ requires a separate topological normalization. Color is an internal axial asymmetry, and the gluon carries transitions between color states.

\subsubsection{Photons and Bosons}

The photon is a massless, helicoidal transverse wave without localized confinement. RefG associates its two helicities and $U(1)$ structure with the transverse mode of the underlying medium. The eight-mode map of the gluons was described in the preceding subsection. In the RefG map, the $W$ and $Z$ bosons are locally confined, massive mediator modes.

RefG describes beta decay as a topological rearrangement: the closed structure changes the distribution of winding, couples to the charged weak channel, and emits a neutral, incompletely confined mode. The physical content of this picture connects the chirality of the weak interaction, energy--momentum balance, and the decay spectrum within one geometrical process.

\subsubsection{Neutrinos and the Higgs Boson}

In this map, the neutrino is an electrically neutral, incompletely confined mode. Incomplete confinement is associated with its small mass and weak coupling to the medium, while flavor oscillation is associated with the phase mixing of internal modes. RefG reads the Higgs boson as a nonrotating, spherical radial pulsation, naturally corresponding to a spin-0 state. Neutrino masses and the mixing matrix, the Higgs mass of $125\,\mathrm{GeV}$, the vacuum expectation value, and the Yukawa couplings are quantitative targets of this geometrical map.

\subsubsection{The Graviton and the Weinberg--Witten Theorem}

The local tensor branch of RefG contains two polarizations, $h_+$ and $h_\times$. The Weinberg--Witten theorem places a strict constraint on Lorentz-covariant conserved tensors and composite massless high-spin states \citep{weinbergwitten1980}. RefG's fundamental medium-based description examines its assumptions separately; the distinguishing features of the low-energy limit are a massless tensor mode, universal coupling, and effective Lorentz symmetry.

In this ontological picture, the graviton is not a separate fundamental substance; it is the quantum mode of a collective tensor displacement of the underlying medium. A purely scalar acoustic oscillation cannot represent the observed gravitational response, because the wave sector needs tensor polarizations and universal coupling. It is also crucial to the Weinberg--Witten constraint that the fundamental level has not been postulated in advance as a Lorentz-covariant field; the recovery of effective Lorentz symmetry in the low-energy oscillon sector is a necessary condition of this picture.

\subsubsection{Dynamics and the Lagrangian}

The mathematical language of this map is the Lagrangian---the dynamical law from which the field equations are obtained and the admissible forms are determined. Even a simple local field model can admit stable knotted solitons \citep{faddeev1997}. RefG connects this mechanism to the particle spectrum, masses, and decay rules.

\subsection{Fundamental Constants: \texorpdfstring{$\alpha$}{alpha} and the Gravitational Hierarchy}

The fine-structure constant, $\alpha\approx1/137$, is one of the well-known mysteries of physics---Feynman regarded it as one of the greatest unexplained numbers of theoretical physics \citep{feynman1985}. Within RefG, $\alpha$ is most plausibly a dimensionless invariant ratio between the electric and magnetic responses of the underlying node network.

The search for $\alpha$ calls for a rigorous self-audit. Simple ratios of Planck-scale and electron-scale masses or frequencies do not give 137; the ratio of the reduced Compton wavelength to the classical electron radius is exactly $1/\alpha$ only because those quantities are defined that way. Nor can the simplicity of the number distinguish 137.036 from its neighbors. Integer steps in phase locking are real physics---the Josephson effect and Shapiro steps demonstrate this \citep{josephson1962,shapiro1963}---but the “second” is itself defined by an atomic oscillon. An absolute frequency in hertz for the underlying medium therefore has no physical meaning, and $\alpha$ must be sought only in dimensionless locking ratios.

The node network of the underlying medium supports electric flux and transverse photon oscillations as two complementary responses of the same microscopic dynamics. The constant $\alpha$ measures their relative strength. When the coupling between nodes strengthens or weakens uniformly, both responses rescale together and their ratio remains unchanged. In this picture, $\alpha$ is not an absolute frequency or energy; it is an intrinsic structural ratio of the network.

The second fundamental question is the \textbf{hierarchy problem}: why is gravity so weak? Between two protons, the electromagnetic force is approximately $10^{36}$ times stronger than gravity; for a pair of electrons, the ratio is of order $10^{42}$. The relation between these large numbers became the historical starting point of Dirac's hypothesis \citep{dirac1937}.

RefG associates the hierarchy with the structural difference between two channels: \textbf{direct phase--transverse coupling} and a \textbf{secondary pressure imprint}.

\begin{itemize}
\item The \textbf{Planck mass}, $M_P=\sqrt{\hbar c/G}$, is the calibration scale for gravitational coupling.
\item The \textbf{electron and proton} are microscopic oscillons with masses many orders of magnitude below the Planck mass.
\end{itemize}

Here gravity is the long-range imprint of the pressure deficit created by an oscillon's energy. The electromagnetic interaction is a direct phase--transverse channel.

For two electrons, the force ratio can be rewritten algebraically as

\begin{equation}
\frac{F_{\rm em}}{F_g}=\alpha\left(\frac{M_P}{m_e}\right)^2.
\label{eq:force-hierarchy}
\end{equation}

This formula displays the hierarchy in terms of the ratio of the Planck mass to the electron mass. RefG associates its microscopic basis with the common origin of $G$, $M_P$, and the oscillon's energy.

In this map, the weakness of gravity reflects the different ontological nature of the two interaction channels. Its quantitative picture unites particle masses, the photon mode, and charge normalization in a single system.

\section{Epilogue: Limits of Knowledge}

\subsection{Why a Single Underlying Medium?}

Let us return, in closing, to the central question: why a single underlying medium? One motivation is the principle of parsimony---a common carrier reduces the number of independent starting assumptions. Yet simplicity alone cannot substantiate a physical theory.

The more important argument concerns the \textbf{origin of regularities}. By treating the electron and the photon as different resonant modes of the same underlying medium, RefG seeks to derive the conservation of energy and charge from the dynamical balance of their common carrier rather than leave these laws as independent starting points. In this language, interactions are couplings between modes of the common medium, and conservation laws are geometrical expressions of their balance.

This argument has a historical background: Eugene Wigner wrote of the “unreasonable effectiveness” of mathematics in the natural sciences \citep{wigner1960}. The ontology of a common carrier offers a physical explanation for this order: shared mathematical regularities reflect a shared dynamical structure.

A common symmetry among many independent fields does not close this question; it merely moves it back one step: where does the unifying rule itself come from, the rule that brings different fields under one law? In the picture of a single underlying medium, this rule is no longer an agreement imposed from outside---it is the dynamics of the common carrier.

This argument also has an empirical form---the problem of precise mutual compatibility. The electrical neutrality of macroscopic matter, the stringent upper bound on the photon mass, and the pattern of fractional quark charges call for a common structural explanation \citep{navas2024}. RefG connects this order to the geometry of modes of a single underlying medium.

One underlying medium gives rise to a rich phenomenology: it can carry pressure, shear, phase, topology, rotation, memory, and resonance at once. This position echoes Philip Anderson's well-known dictum, \textit{More is different}: each new level of organization has regularities of its own \citep{anderson1972}. RefG reads the transitions between levels as the dynamics of this multichannel underlying medium.

RefG reads the Universe as a finite, boundaryless global resonator. The discrete spectrum discussed in Chapter IV gives this picture a micro--macro spectral foundation, while its physical imprint is read in cosmological modes and observable topology.

Three principal ontological views concern us in this article: an infinite underlying medium; a finite Universe viewed, as it were, from outside; and a finite underlying medium that measures itself only from within, through its own modes. The second view calls for a new level of external observer and generates a regress; the first takes actual infinity as an initial given. RefG chooses the third picture---a finite, boundaryless, self-reflexive resonator. Its physical meaning is read in its own spectrum and in the imprint left within the horizon. Here the Universe measures itself through its own resonances.

\subsection{Three Operational Infinities}

The Universe is always measured by something that is itself part of the Universe: our clocks, rulers, detectors, and bodies are all oscillon modes of the universal underlying medium. This concept of a “self-measuring Universe” is related to Wheeler's \textit{It from Bit} program \citep{wheeler1990}. For Wheeler, information is primary; in RefG, the underlying medium as substance is primary, and the informational imprint is a result of its self-distinction.

In what follows, “infinity” has an operational meaning: it denotes a limit or extension of a measurement procedure and does not presuppose that the Universe is in fact infinite in size, age, or divisibility.

Here Gödel's theorem plays the role of a logical analogy. In any sufficiently strong, consistent formal system, there are propositions that cannot be decided from within that system \citep{godel1931}. An analogous question arises in a self-measuring Universe: how completely can one part of the system measure the boundaries of the whole? From this intuition emerge the “three operational infinities”:

\textbf{Temporal infinity:} Before the first resonant clocks, operational time was undefined. “13.8 billion years” is metric time calibrated by present-day clocks. A process interpretation expresses the internal history of the early underlying medium as a larger budget of intrinsic cycles, while metric cosmological time retains the same 13.8-billion-year calibration.

\textbf{Spatial infinity:} We do not measure “pure” emptiness; we measure only the relations between distinguishable nodes. The cosmological horizon is therefore the operational boundary of our measurable network, not the physical end of the underlying medium itself. The region beyond the horizon remains outside direct observation.

An indirect signature is nevertheless possible: if the size of the global resonator is of the order of the horizon, its longest allowed mode will leave an imprint on low-multipole power and cosmological topology. This relation is defined through a spectral imprint measurable within our horizon.

\textbf{Scale infinity:} In RefG, the Planck length denotes not the absolute floor of the underlying medium but a natural crossover scale built from $c$, $\hbar$, and $G$, where quantum and gravitational effects become important together. In the history of physics, the “smallest” known scale has changed more than once---from the atom to the nucleon, and from the nucleon to the quark. Below the Planck scale there may be a deeper dynamical layer, perhaps with its own effective constants, or no new physical structure may be found in that direction at all. Larger organized structures may also exist above our scale. RefG bounds both directions by measurable imprints.

In all three cases, infinity arises between the measurement procedure and the demand for a complete description. Even a finite object can unfold without bound in a chosen representation: stereographic projection maps a finite sphere onto an infinite plane. Mandelbrot's “coastline paradox” expresses a related idea, in which the measured length increases as the measuring scale decreases \citep{mandelbrot1982}. RefG reads these examples as operational properties of description and measurement scale.

\subsection{Methodological Principle}

The methodological principle of this article is clear: an intuitive physical picture must become a rigorous, calculable, and falsifiable formalism. The analysis of a particle oscillon follows this sequence:

The text clearly distinguishes a working hypothesis from a computational result: a result is always tied to a specific code, the input data used, and the output obtained.

\begin{enumerate}
\item \textbf{Experimental basis}---record the mass, charge, spin, lifetime, and other observables.
\item \textbf{Mode identification}---select the corresponding resonant class of the underlying medium.
\item \textbf{Initial geometrical form (ansatz)}---describe the oscillon's form, topology, and boundary conditions precisely.
\item \textbf{Dynamical law}---derive the action, field equations, and source.
\item \textbf{Symmetries and conservation laws}---determine the charge, spin, color, and zero modes.
\item \textbf{Existence and stability}---test the finite-energy solution, matching, and perturbation spectrum.
\item \textbf{Parameter closure}---determine every coefficient independently of the data we seek to predict.
\item \textbf{New prediction}---calculate an independent observable and its uncertainty.
\item \textbf{Data test}---compare the result with observations using a predefined falsification criterion.
\end{enumerate}

The aim of this sequence is to derive the gyromagnetic factor from oscillon dynamics and to turn the formula “topological locking = charge” into a measurable, independently testable law.

The same rule governs the geometrical origin of hypercharge and isospin, a band-gap-like mechanism of confinement, the evolution of the background temperature, the emergence of one effective metric seen by every material mode, and separate tests for a preferred frame or spatial anisotropy. If the intuition cannot be turned into a calculable formalism, the intuition must be revised.

\subsection{Where This Program Can Die}

The list of working tasks states what we must do; the falsification map states what nature can remove. The outcomes divide into three levels: the entire framework may collapse, a particular physical sector may fail, or only a specific geometrical reading may fail. This distinction matters, because the failure of one branch does not automatically nullify every other idea, whereas a violation of the common foundational principle cannot be repaired by reworking an individual branch.

\textbf{The entire framework is rejected} if any of the following conditions is met:

\begin{enumerate}
\item Effective Lorentz symmetry and the isotropy of the speed of light fail to emerge in the oscillon sector. Optical resonators constrain anisotropy at the $10^{-17}$ level \citep{herrmann2009}, while the gravitational and electromagnetic signals from GW170817 tie the tensor-wave speed to the speed of light at approximately the $10^{-15}$ level \citep{abbott2017}.
\item The inertial and gravitational responses tied to the same total energy leave a composition-dependent residual. The MICROSCOPE experiment constrains such a difference at the $10^{-15}$ level \citep{touboul2022}.
\item Process time has to be added to the expansion history, the CMB, BBN, photon redshift, or the $(1+z)$ time dilation of supernovae. The process language cannot replace the metric calculation of these channels \citep{blondin2008}.
\item Dimensionless ratios change under a common transposition. Quasar spectra test $\Delta\alpha/\alpha$ at the part-per-million level \citep{kotus2017}; such a change rejects the principle of the local collective rhythm.
\end{enumerate}

\textbf{A particular sector is rejected} if:

\begin{enumerate}
\item The source--flux normalization of gravitational waves fails to recover the quadrupolar energy loss tested by the binary pulsar to a precision of $10^{-3}$--$10^{-4}$, or a scalar pulsational polarization exceeds the observational bound \citep{kramer2021,ligovirgokagra2022}.
\item The galactic-vortex map fails to describe SPARC/RAR, the mass budget of clusters, and the approximately $8\sigma$ separation between the lensing centers and the X-ray plasma in the Bullet Cluster \citep{clowe2006}. This result kills the vortex sector that replaces dark matter, not the refractive weak-field chain.
\item With the same metric, matter content, recombination, and initial conditions, the identity of the primary CMB sources is violated, or the combined vortical--nodal sources replacing dark matter fail to recover the likelihood structure of the CMB/LSS and the BAO standard ruler \citep{eisenstein2005}. The first result rejects the consistency bridge; the second rejects the cosmological replacement without dark matter.
\item A specific compact branch fails to match the shadow size, oscillation spectrum, and bounds on echo signals. Such a result rejects precisely the internal geometry whose external phenomenology was tested \citep{eventhorizon2019,cardoso2019}.
\end{enumerate}

\textbf{A particular geometrical reading is rejected} if:

\begin{enumerate}
\item An extended oscillon profile of the electron fails to recover the observed pointlike form factor.
\item The Koide/$C_3$ map fails to find a definition of mass and a radiative bridge that bring the exact branch $(A_K,\theta)=(\sqrt2,2/9)$ into agreement with the measured lepton masses. A fourth sequential charged lepton rejects the strict, exhaustive three-mode version of this map.
\item Once the effective mass and the coherence window have been fixed independently, circulation in the galactic flow is not quantized in units of $h/m$. This result rejects the particular reading in terms of a discrete vortical medium, not the existence of a continuous vortex.
\item The dispersion of Kelvin waves derived from the oscillon action fails a specified comparison with classical and quantum vortices.
\item The Möbius geometry fails to produce the Dirac value $g=2$, or its quantum completion fails to recover the radiative corrections of the electron. At the same level, failure to recover the antisymmetry of Fock space nullifies the geometrical mechanism for the Pauli principle.
\end{enumerate}

Each item has a specific measurement or a decisive calculation. This is why the list marks the boundary of the research program: a negative answer from nature removes the corresponding idea from the program rather than merely changing its wording.

\section*{Conclusion}
\addcontentsline{toc}{section}{Conclusion}

This article has built a single, continuous conceptual chain grounded in the principles of Refractive Gravity (RefG). Its ontological picture of the Universe is this: the spontaneous self-differentiation, or self-distinction, of one underlying medium gives rise to the onset of measurability---the operational seeds of space, time, and mass. A locally confined resonant form of the medium, the \textbf{oscillon}, creates a pressure deficit in the medium that an external observer reads as effective mass. Motion within the effective geometry altered by this deficit is gravity---the universal refractive response to the state of the medium.

This unified pressure chain unfolds across three principal scales:

\begin{itemize}
\item at the \textbf{local level}---in compact objects, where internal state and external readout are clearly separated;
\item at the \textbf{galactic level}---in the large-scale vortical flow of the underlying medium, which RefG connects to MOND phenomenology;
\item at the \textbf{cosmological level}---in the historical, asymptotic relaxation of the state of the underlying medium.
\end{itemize}

The allowed spectrum of the underlying medium and mutual population locking connect the sparse, discrete particle spectrum to a single mechanism, while the topology of phase and orientation gives quantum numbers a geometrical language. Along the same line, the fine-structure constant $\alpha$, interpreted as an invariant ratio between the electric and magnetic responses of the node network, connects the oscillon map to the dimensionless structure of the electromagnetic interaction.

The article rests on the weak-field regime---the domain in which gravity has been tested most precisely. On the selected 1PN branch, the refractive language of RefG agrees with the standard description of general relativity. From this foundation, the research program extends to particle, galactic, and cosmological scales, where the physical map organizes completed calculations and predefined observational tests into one system.

The value of this picture lies not only in its breadth of unification. Every important connection leaves nature a specific test: Lorentz symmetry, equivalence, wave-energy balance, galactic dynamics, the particle form factor, and topological numbers. Measurement therefore assembles this puzzle; a piece that does not fit nature's answer must be changed.

From a single ontological foundation, Refractive Gravity develops a research program whose aim is a unified, calculable, and empirically testable description of gravity, matter, and measurability. The present text lays the foundation of this physical paradigm and opens the way for further research.

\section*{Current Status of the Research Program}
\addcontentsline{toc}{section}{Current Status of the Research Program}

Current status information is gathered here in one place so that the main text can preserve the continuity of the physical narrative.

\begin{enumerate}
\item \textbf{The underlying medium and the origin of measurability}---the ontological core has been formulated; operational distinguishability defines the starting point from which space, time, and mass emerge.

\item \textbf{Inter-node separation and minimum distinguishability}---the idea of operational distinguishability has been formulated; a threshold of the order of the Planck scale is a physical candidate and has not been derived from the action.

\item \textbf{The local collective rhythm}---common-transposition invariance and toy models of mutual locking have been obtained; calculation of the actual stable oscillon set remains open.

\item \textbf{Oscillon deficit, size, time, and mass}---the metric filters and the first-order biconformal branch have been obtained symbolically; a finite-energy nonlinear oscillon profile and the complete mass--readout map remain open.

\item \textbf{Weak-field gravity}---on the selected 1PN branch, $\gamma=1$, $\beta=1$, and the listed observables agree with the standard results; higher orders, a complete PPN audit of preferred-frame effects, the description of all material modes by one effective metric, and a full likelihood-based test against the data remain open.

\item \textbf{Inertia and gravitational mass}---the Noether identity has been obtained conditionally for a finite-energy, Lorentz-covariant localized solution satisfying the Laue boundary condition; a direct derivation of the zero mode remains open.

\item \textbf{Gravitational waves and local unobservability}---in the local Minkowski limit, the tensor propagation speed and the transverse--traceless detector response have been obtained; cosmological propagation, source--flux normalization, the complete waveform, and scalar polarization remain open.

\item \textbf{Compact objects}---a static, conditional map and the finite-core matching of the $C^2$ ansatz have been obtained; in one specific set of exterior-layer and pressure-reversal criteria, the admissible domains do not overlap. The physical equation of state, collapse, rotation, and full stability remain open.

\item \textbf{The galactic vortex and MOND/Tully--Fisher}---the transition algebra and the baryonic Tully--Fisher relation (BTFR) follow conditionally after the assumption $a_0=cH/(2\pi)$ and vortex closure; derivation from the action, SPARC/RAR, the coefficients of the three source channels in clusters, and complete modeling of clusters remain open.

\item \textbf{Cosmological relaxation and dark energy}---the Friedmann--Lemaître--Robertson--Walker (FLRW) stress algebra and a conditional $\Lambda$-like diagnostic have been obtained; pressure relaxation is an interpretive bridge. Perturbations and numerical fits to BBN, CMB, Planck, and BAO remain open.

\item \textbf{Two-level time and quantum randomness}---the interpretive map is compatible with the facts of Bell tests and delayed-choice experiments; the Born distribution, the no-signaling condition, and the complete dynamics of measurement have not been derived.

\item \textbf{Particles as oscillon modes}---the geometrical placement map and initial spectral scheme have been formulated as candidates; a finite-energy oscillon, its form factor, masses, couplings, and decays remain open.

\item \textbf{The Koide/$C_3$/$2/9$ structure}---the $C_3$ algebra gives $K=2/3$ exactly; with $(A_K,\theta)=(\sqrt2,2/9)$ and the electron pole mass as the anchor, the result misses the muon mass by $451.7\sigma$ and the tau mass by $0.61\sigma$. The choice $h=2$, the relation $m\propto\nu^2$, derivation from the action, and radiative protection remain open.

\item \textbf{Spin, Pauli exclusion, and fermionic statistics}---the Möbius double-cover geometry provides a candidate mechanism; a complete derivation of Fock space and of the spin--statistics relation of quantum field theory remains open.

\item \textbf{Charge, the photon, and the electromagnetic channel}---subtheorems establishing local-frame $U(1)$ symmetry, two helicities, a screened far field, and a conserved phase current have been obtained; canonical charge normalization, Maxwell's equations, Gauss's law, the Ward identities, quantum electrodynamics (QED), the Dirac sector, and the exact pressure-deficit channel through which free radiation sources gravity remain open.

\item \textbf{Color, quarks, and gluons}---the three-color, eight-mode construction provides only a partially compatible map; the structure constants of $SU(3)$, interaction vertices, scale evolution, the confinement potential, and fractional charge remain open.

\item \textbf{The electroweak map and the Higgs}---this is a geometrical placement map; $SU(2)_L$ chirality, a $W/Z$ mass matrix that preserves a massless photon, CKM/PMNS mixing, the neutrino mass ordering, Yukawa couplings, and the vacuum expectation value have not been derived.

\item \textbf{The fine-structure constant $\alpha$}---a node-based interpretation and a target-independent spectral measure have been obtained: $\alpha$ is read from the relative strengths of the electric-flux and photon responses, while the overall scale cancels. A compact microscopic action for the physical charge generator $Q$, a unique numerical ratio, the complete charged spectrum, and evolution to the Thomson limit remain open. The earlier boundary formula remains a target-dependent conditional result.

\item \textbf{The gravitational hierarchy}---the force ratio can be written algebraically as $F_{\rm em}/F_g=\alpha(M_P/m)^2$; the origins of $G$, $M_P$, the oscillon's microscopic energy scale, and the mass hierarchy remain open.

\item \textbf{A finite self-measuring Universe}---the global resonator is a cosmological hypothesis; low-multipole and topological signatures are candidates, and no unique likelihood fit to the data has been obtained.
\end{enumerate}

\phantomsection
\addcontentsline{toc}{section}{References}
\bibliographystyle{plainnat}
\bibliography{RefG_EN}

\end{document}
